\documentclass[11pt,a4paper]{article}
\usepackage[utf8]{inputenc}
\usepackage[T1]{fontenc}
\usepackage[french]{babel}
\usepackage{geometry}
\usepackage{booktabs}
\usepackage{longtable}
\usepackage{array}
\usepackage{xcolor}
\usepackage{colortbl}
\usepackage{graphicx}
\usepackage{fancyhdr}
\usepackage{titlesec}
\usepackage{tikz}
\usepackage{enumitem}
\usepackage{hyperref}
\usepackage{multicol}
\usepackage{listings}
\usepackage{tcolorbox}
\usepackage{amssymb}  
\usepackage{float}  % <--- AJOUTER CETTE LIGNE



\geometry{margin=2cm}
\pagestyle{fancy}
\fancyhf{}
\rhead{Atelier Machine Learning}
\lhead{Analyse Comportementale Clientèle}
\rfoot{Page \thepage}

% Couleurs personnalisées
\definecolor{headerblue}{RGB}{41, 65, 114}
\definecolor{lightgray}{RGB}{245, 245, 245}
\definecolor{accentorange}{RGB}{230, 126, 34}
\definecolor{codegreen}{RGB}{0,128,0}
\definecolor{codegray}{RGB}{128,128,128}
\definecolor{codepurple}{RGB}{128,0,128}
\definecolor{backcolour}{RGB}{248,248,248}
\definecolor{warningred}{RGB}{231, 76, 60}
\definecolor{successgreen}{RGB}{46, 204, 113}

\lstdefinestyle{mystyle}{
    backgroundcolor=\color{backcolour},   
    commentstyle=\color{codegreen},
    keywordstyle=\color{codepurple},
    numberstyle=\tiny\color{codegray},
    stringstyle=\color{codegreen},
    basicstyle=\ttfamily\footnotesize,
    breakatwhitespace=false,         
    breaklines=true,                 
    captionpos=b,                    
    keepspaces=true,                 
    numbers=left,                    
    numbersep=5pt,                  
    showspaces=false,                
    showstringspaces=false,
    showtabs=false,                  
    tabsize=2
}
\lstset{style=mystyle}


\tcbset{
    colback=backcolour,
    colframe=headerblue,
    fonttitle=\bfseries,
    boxrule=1pt,
    arc=3mm
}


\titleformat{\section}
{\color{headerblue}\normalfont\Large\bfseries}
{\thesection}{1em}{}[{\titlerule[2pt]}]

\titleformat{\subsection}
{\color{accentorange}\normalfont\large\bfseries}
{\thesubsection}{1em}{}

% Définition de colonnes personnalisées
\newcolumntype{C}[1]{>{\centering\arraybackslash}p{#1}}
\newcolumntype{L}[1]{>{\raggedright\arraybackslash}p{#1}}

\begin{document}


\begin{titlepage}
\begin{tikzpicture}[remember picture,overlay]
    \fill[headerblue] (current page.north west) rectangle ([yshift=-5cm]current page.north east);
    \node[anchor=north west, xshift=2cm, yshift=-2cm, text=white] at (current page.north west) 
        {\Huge\textbf{Atelier Machine Learning}};
    \node[anchor=north west, xshift=2cm, yshift=-3.5cm, text=white] at (current page.north west) 
        {\Large\textit{Analyse Comportementale Clientèle Retail}};
\end{tikzpicture}

\vspace{6cm}

\begin{center}
{\LARGE\textbf{COMPTE RENDU}}\par
\vspace{1cm}
{\large Atelier Pratique : E-commerce de Cadeaux}\par
\vspace{2cm}
{\large \textbf{Préparé par Ranim Bouguila}}
\vspace{3cm}

\begin{tikzpicture}
    \draw[headerblue, line width=2pt] (0,0) -- (10,0);
\end{tikzpicture}

\vspace{2cm}



\vspace{5cm}

{\textbf{Année Univeristaire : 2025-2026}}

\end{center}
\vfill
\end{titlepage}

% ================================
% TABLE DES MATIÈRES
% ================================
\tableofcontents
\newpage

% ================================
% LISTE DES FIGURES
% ================================
\listoffigures
\newpage

% ================================
% LISTE DES TABLEAUX
% ================================
\listoftables
\newpage

\section{Introduction}
Ce projet vise à construire une chaîne complète de Machine Learning pour analyser le comportement des clients d'un e-commerce de cadeaux. Deux objectifs principaux :
\begin{itemize}
	\item \textbf{Segmenter} la clientèle (VIP, clients dormants, nouveaux)
	\item \textbf{Prédire} le churn (départ des clients)
\end{itemize}
Ce rapport détaille la première phase : le \textbf{prétraitement des données}.

\section{Architecture du Pipeline}
Le pipeline de prétraitement se décompose en 6 étapes successives :
\begin{figure}[H]
	\centering
	\includegraphics[width=0.5\textwidth]{images/pipeline.png}
	\caption{Pipeline complet de prétraitement}
	\label{fig:pipeline}
\end{figure}
\subsection{Variables RFM (comportement d'achat)}
\begin{itemize}
	\item \texttt{Recency} : nombre de jours écoulés depuis le dernier achat (0-400 jours)
	\item \texttt{Frequency} : nombre de commandes distinctes passées (1-248)
	\item \texttt{MonetaryTotal} : dépense totale en livres sterling, peut être négative (remboursements), de -4287£ à +279489£
\end{itemize}
\subsection{Variables Produit}
\begin{itemize}
	\item \texttt{TotalQuantity}, \texttt{UniqueProducts}, \texttt{UniqueDescriptions}
	\item \texttt{ReturnRatio}, \texttt{CancelledTransactions}, \texttt{NegativeQuantityCount}
\end{itemize}
\subsection{Variables Client}
\begin{itemize}
	\item \texttt{Age} : 30\% de valeurs manquantes (1311 NaN sur 4372)
	\item \texttt{SupportTicketsCount} et \texttt{SatisfactionScore} : contiennent des valeurs sentinelles (-1, 999, 0, 99) encodant des erreurs système
	\item \texttt{Gender} (M/F/Unknown), \texttt{Region} (13 régions), \texttt{Country} (37 pays)
	\item \texttt{Churn} : variable cible binaire (0=fidèle, 1=parti), taux=33.3\%
\end{itemize}
\section{Étapes du Prétraitement}

\subsection{Étape 1 — Chargement et Analyse Initiale}
Diagnostic complet du dataset : types de données, statistiques descriptives, valeurs manquantes par colonne, doublons.
\textbf{Résultat} : 0 doublon détecté. 2 colonnes avec NaN :
\begin{itemize}
	\item \texttt{Age} : 1311 valeurs manquantes (30\%)
	\item \texttt{AvgDaysBetweenPurchases} : 79 valeurs manquantes (1.8\%)
\end{itemize}
\subsection{Étape 2 — Nettoyage de Base}
\begin{itemize}
	\item Suppression des doublons : 0 supprimé
	\item Imputation de \texttt{Age} par la médiane (49.0 ans) — la médiane est préférée à la moyenne car robuste aux valeurs extrêmes
	\item Imputation de \texttt{AvgDaysBetweenPurchases} par la médiane
\end{itemize}
\subsection{Étape 3 — Valeurs Aberrantes et Sentinelles}
Deux types de problèmes distincts sont traités :

\textbf{Type 1 — Valeurs sentinelles (codes d'erreur déguisés en nombres)} :
\begin{itemize}
	\item \texttt{SupportTicketsCount} : 130 valeurs -1 et 999 remplacées par NaN puis imputées par la médiane (2.0 tickets)
	\item \texttt{SatisfactionScore} : 349 valeurs -1, 0 et 99 remplacées par NaN puis imputées par la médiane (3.0/5)
\end{itemize}
\textbf{Type 2 — Vraies valeurs aberrantes statistiques} :
	Méthode IQR (Interquartile Range) appliquée sur \texttt{SupportTicketsCount} : bornes calculées $[Q1-1.5\times IQR, Q3+1.5\times IQR] = [-2.00, 6.00]$. 23 valeurs extrêmes ramenées aux bornes (capping).
\begin{figure}[H]
	\centering
	\includegraphics[width=0.9\textwidth]{images/iqr_method.png}
	\caption{Traitement des valeurs aberrantes - Méthode IQR}
	\label{fig:iqr}
\end{figure}
\subsection{Étape 4 — Formats Inconsistants et Feature Engineering}
\textbf{RegistrationDate} : formats hétérogènes (DD/MM/YYYY, YYYY-MM-DD, etc.) convertis en datetime uniforme. 4 nouvelles features extraites : \texttt{RegYear}, \texttt{RegMonth}, \texttt{RegDay}, \texttt{RegWeekday}. La colonne date originale est ensuite supprimée.

\textbf{LastLoginIP} : adresse IP brute transformée en 5 features exploitables : \texttt{IP1}, \texttt{IP2}, \texttt{IP3}, \texttt{IP4} (4 octets numériques) et \texttt{IP\_Type} (Private/Public selon les plages RFC 1918 : 10.x.x.x, 172.x.x.x, 192.x.x.x). Résultat : 4321 IPs publiques, 51 IPs privées.

\textbf{NewsletterSubscribed} : supprimée car variance nulle (toujours "Yes") — aucune information prédictive.
\subsection{Étape 5 — Encodage des Variables Catégorielles}
3 stratégies selon la nature de chaque colonne :

\textbf{Encodage Ordinal (6 colonnes avec ordre logique)} :
\begin{itemize}
	\item \texttt{SpendingCategory} : Low=0, Medium=1, High=2, VIP=3
	\item \texttt{AgeCategory} : Inconnu=0, 18-24=1, ..., 65+=6
	\item \texttt{LoyaltyLevel} : Nouveau=0, Jeune=1, Établi=2, Ancien=3
	\item \texttt{ChurnRiskCategory} : Faible=0, Moyen=1, Élevé=2, Critique=3
	\item \texttt{BasketSizeCategory} : Petit=0, Moyen=1, Grand=2
	\item \texttt{PreferredTimeOfDay} : Nuit=0, Matin=1, Midi=2, Après-midi=3, Soir=4
\textbf{Encodage One-Hot (8 colonnes sans ordre)} :
\texttt{Gender}, \texttt{Region}, \texttt{AccountStatus}, \texttt{CustomerType}, \texttt{FavoriteSeason}, \texttt{WeekendPreference}, \texttt{ProductDiversity}, \texttt{IP\_Type}. \texttt{drop\_first=True} pour éviter la multicolinéarité. Résultat : 8 colonnes → 29 nouvelles colonnes binaires.
\textbf{Encodage par Fréquence (Country — 37 valeurs uniques)} :
Chaque pays remplacé par sa proportion dans le dataset. Ex : United Kingdom → 0.904, France → 0.020, Germany → 0.022.
\subsection{Étape 6 — Split Train/Test et Normalisation}
Split stratifié 80/20 avec \texttt{random\_state=42} pour la reproductibilité. \texttt{stratify=y} garantit que les proportions Churn/Non-Churn sont conservées dans les deux sous-ensembles.

\begin{figure}[H]
	\centering
	\includegraphics[width=1\textwidth]{images/train_test_split.png}
	\caption{éviter le data leakage}
	\label{fig:split}
\end{figure}

\textbf{StandardScaler} appliqué uniquement sur X\_train (fit\_transform), puis transform sur X\_test — règle absolue pour éviter le data leakage. La variable cible y n'est jamais normalisée.
\end{itemize}
\section{Résultats du Prétraitement}

\begin{table}[H]
	\centering
	\caption{Résumé des transformations}
	\label{tab:transformations}
	\begin{tabular}{|l|c|}
		\hline
		\textbf{Opération} & \textbf{Résultat} \\
		\hline
		Doublons & 0 trouvés \\
		\hline
		Imputation Age & 1311 valeurs → médiane 49.0 \\
		\hline
		AvgDaysBetweenPurchases & 79 valeurs → médiane \\
		\hline
		SupportTicketsCount & 130 sentinelles + 23 outliers corrigés \\
		\hline
		SatisfactionScore & 349 sentinelles corrigées \\
		\hline
		RegistrationDate & 4 features extraites (Year, Month, Day, Weekday) \\
		\hline
		NewsletterSubscribed & Supprimée (variance nulle) \\
		\hline
		LastLoginIP & $\rightarrow$ IP1, IP2, IP3, IP4, IP\_Type \\
		\hline
		Encodage ordinal & 6 colonnes encodées \\
		\hline
		Encodage One-Hot & 8 colonnes $\rightarrow$ 29 nouvelles colonnes \\
		\hline
		Encodage fréquence & Country $\rightarrow$ Country\_freq \\
		\hline
		\textbf{Dimensions finales} & \textbf{4372 lignes × 81 colonnes} \\
		\hline
		\end{tabular}
	\end{table}
\begin{table}[H]
	\centering
	\caption{Dataset final prêt pour la modélisation}
	\label{tab:final_data}
	\begin{tabular}{|l|c|c|}
		\hline
		\textbf{Fichier} & \textbf{Lignes} & \textbf{Colonnes} \\
		\hline
		X\_train.csv & 3497 & 79 \\
		X\_test.csv & 875 & 79 \\
		y\_train.csv & 3497 & 1 \\
		y\_test.csv & 875 & 1 \\
		\hline
	\end{tabular}
\end{table}
\subsection{Observations Importantes}
\begin{itemize}
	\item \textbf{Taux de churn} : 33.3\% (2918 non-churn / 1454 churn) → déséquilibre modéré, traité par \texttt{class\_weight='balanced'} lors de la modélisation
	\item \textbf{Région} : 90.5\% UK → déséquilibre géographique important
	\item \textbf{WeekendPreference} : 71.5\% "Inconnu" → feature peu informative
\end{itemize}
\section{Clustering K-Means}

\subsection{Objectif}
Segmenter les clients en groupes homogènes selon leur comportement d'achat sans utiliser la variable cible \texttt{Churn} (apprentissage non supervisé).

\subsection{Features Utilisées}
\texttt{Recency}, \texttt{Frequency}, \texttt{MonetaryTotal} (déjà normalisées par StandardScaler lors du prétraitement).

\subsection{Sélection du Nombre de Clusters — Méthode Elbow}
On entraîne K-Means pour k de 2 à 9 et on mesure l'inertie intra-cluster. L'inertie décroît rapidement jusqu'à k=4 puis se stabilise progressivement. Le "coude" identifié à k=4 correspond aux 4 segments RFM du dataset.
\begin{table}[H]
	\centering
	\caption{Inertie par nombre de clusters}
	\label{tab:inertia}
	\begin{tabular}{|c|c|}
		\hline
		\textbf{k} & \textbf{Inertie} \\
		\hline
		2 & 7129 \\
		3 & 4254 \\
		4 & 3228 \\
		5 & 2428 \\
		6 & 1941 \\
		\hline
	\end{tabular}
\end{table}
\begin{figure}[H]
	\centering
	\includegraphics[width=0.8\textwidth]{images/elbow_curve.png}
	\caption{Courbe d'Elbow — sélection du k optimal (coude à k=4)}
	\label{fig:elbow}
\end{figure}
\subsection{Résultats K-Means avec k=4}
\textbf{Silhouette Score} = 0.6041 (supérieur à 0.5 → bon clustering)

\begin{table}[H]
	\centering
	\caption{Profil des 4 clusters clients}
	\label{tab:clusters}
	\begin{tabular}{|c|c|c|c|c|c|}
		\hline
		\textbf{Cluster} & \textbf{Taille} & \textbf{Recency} & \textbf{Frequency} & \textbf{Monetary} & \textbf{Interprétation} \\
		\hline
		0 & 877 & +1.54 & -0.34 & -0.19 & Clients Dormants \\
		1 & 2571 & -0.51 & +0.01 & -0.04 & Clients Réguliers \\
		2 & 44 & -0.86 & +5.19 & +3.54 & Clients Champions \\
		3 & 5 & -0.83 & +8.80 & +21.12 & Clients VIP Absolus \\
		\hline
	\end{tabular}
\end{table}
\subsection{Interprétation des Clusters}
\begin{itemize}
	\item \textbf{Cluster 0 (Dormants, 20\%)} : Recency élevée → n'ont pas acheté depuis longtemps. Fréquence et montant faibles. Risque de churn élevé.
	\item \textbf{Cluster 1 (Réguliers, 59\%)} : comportement moyen équilibré. Cœur de clientèle stable.
	\item \textbf{Cluster 2 (Champions, 1\%)} : achats très récents, très fréquents et montants élevés. Clients à fidéliser en priorité.
	\item \textbf{Cluster 3 (VIP, 0.1\%)} : comportement extrême, dépenses 21× supérieures à la moyenne. Clients ultra-premium.
\end{itemize}
\begin{figure}[H]
	\centering
	\includegraphics[width=0.8\textwidth]{images/kmeans_clusters.png}
	\caption{Segmentation K-Means — Recency vs MonetaryTotal (k=4)}
	\label{fig:clusters}
\end{figure}
\section{Identification et Traitement du Data Leakage}

Lors de l'entraînement initial, les trois classificateurs atteignaient une accuracy de 100\%, révélant un problème de \textbf{data leakage}.

\subsection{Type 1 — Leakage Direct (colonnes calculées depuis Churn)}
Certaines colonnes ont été construites algorithmiquement À PARTIR de la variable \texttt{Churn} lors de la génération du dataset synthétique.
\begin{table}[H]
	\centering
	\caption{Colonnes supprimées pour leakage direct}
	\label{tab:leakage_direct}
	\begin{tabular}{|l|c|l|}
		\hline
		\textbf{Colonne} & \textbf{Corrélation} & \textbf{Raison de suppression} \\
		\hline
		ChurnRiskCategory & 0.878 & Encodage du risque de churn \\
		CustomerType\_Perdu & 0.699 & "Perdu" = churné par définition \\
		RFMSegment\_Dormants & 0.573 & Dormant = inactif = churné \\
		LoyaltyLevel & 0.429 & Niveau de fidélité = inverse du churn \\
		RFMSegment\_Fidèles & 0.300 & Fidèle = non churné par définition \\
		RFMSegment\_Potentiels & 0.299 & Segment à risque de churn \\
		CustomerType\_Nouveau & 0.232 & Corrélé fortement au churn \\
		\hline
	\end{tabular}
\end{table}
\subsection{Type 2 — Leakage Indirect (corrélations artificielles)}
La variable \texttt{Churn} a été générée algorithmiquement à partir de features temporelles, créant des corrélations irréalistes.

\begin{table}[H]
	\centering
	\caption{Colonnes supprimées pour corrélations artificielles}
	\label{tab:leakage_indirect}
	\begin{tabular}{|l|c|l|}
		\hline
		\textbf{Feature} & \textbf{Corrélation} & \textbf{Décision} \\
		\hline
		Recency & 0.858 & Supprimée pour classification \\
		CustomerTenureDays & 0.453 & Supprimée \\
		PreferredMonth & 0.427 & Supprimée \\
		SpendingCategory & 0.384 & Supprimée \\
		FavoriteSeason\_Printemps & 0.302 & Supprimée \\
		Frequency & 0.217 & Conservée (corrélation acceptable) \\
		UniqueInvoices & 0.217 & Conservée \\
		\hline
	\end{tabular}
\end{table}
Au total : \textbf{12 colonnes supprimées} → X passe de 79 à 67 features.
\subsection{Impact de la Correction}
\begin{table}[H]
	\centering
	\caption{Comparaison avant/après suppression du leakage}
	\label{tab:leakage_impact}
	\begin{tabular}{|l|c|c|}
		\hline
		\textbf{Modèle} & \textbf{Avant suppression} & \textbf{Après suppression} \\
		\hline
		Logistic Regression & 100\% & 98.29\% \\
		Random Forest & 100\% & 98.51\% \\
		XGBoost & 100\% & 99.77\% \\
		\hline
	\end{tabular}
\end{table}
\begin{tcolorbox}[colback=warningred!10, colframe=warningred, title=\textbf{Conclusion Méthodologique}]
	Les scores restent élevés (98-99\%) car le dataset est synthétique et pédagogique — le Churn y a été généré par une formule déterministe.
	
	\textbf{Dans un dataset réel de e-commerce :}
	\begin{itemize}
		\item Corrélations features/Churn seraient de 0.15 à 0.35
		\item Performances attendues : 75\% à 88\% d'accuracy
		\item Présence de bruit, cas ambigus, comportements imprévisibles
	\end{itemize}
	
	\textbf{Ce problème illustre l'importance de :}
	\begin{enumerate}
		\item Analyser les corrélations features/target avant entraînement
		\item Identifier les colonnes construites depuis la target
		\item Distinguer performances sur données synthétiques vs réelles
	\end{enumerate}
\end{tcolorbox}

\section{Classification — Prédiction du Churn}

\subsection{Objectif}
Prédire si un client va partir (Churn=1) ou rester (Churn=0). Features : 67 colonnes après suppression du leakage.

\subsection{Gestion du Déséquilibre}
\begin{itemize}
	\item Logistic Regression et Random Forest : \texttt{class\_weight='balanced'}

\end{itemize}
subsection{Résultats}

\begin{table}[H]
	\centering
	\caption{Performances des modèles de classification}
	\label{tab:classification}
	\begin{tabular}{|l|c|c|c|c|}
		\hline
		\textbf{Modèle} & \textbf{Accuracy} & \textbf{F1-Score} & \textbf{FP} & \textbf{FN} \\
		\hline
		Logistic Regression & 98.29\% & 0.9828 & 5 & 10 \\
		Random Forest & 98.51\% & 0.9851 & 0 & 13 \\
		\textbf{XGBoost} & \textbf{99.77\%} & \textbf{0.9977} & \textbf{0} & \textbf{2} \\
		\hline
	\end{tabular}
\end{table}
\begin{figure}[H]
	\centering
	\includegraphics[width=0.5\textwidth]{images/confusion_logistic_regression.png}
	\includegraphics[width=0.5\textwidth]{images/confusion_random_forest.png}
	\includegraphics[width=0.5\textwidth]{images/confusion_xgboost.png}
	\caption{Matrices de confusion des trois modèles}
	\label{fig:confusion}
\end{figure}

\subsection{Importance des Features (XGBoost)}
\begin{figure}[H]
	\centering
	\includegraphics[width=0.8\textwidth]{images/feature_importance_classification.png}
	\caption{Top 10 des features les plus importantes pour la prédiction du churn}
	\label{fig:feature_importance}
\end{figure}
\begin{table}[H]
	\centering
	\caption{Top 5 features importantes (XGBoost)}
	\label{tab:features}
	\begin{tabular}{|c|c|}
		\hline
		\textbf{Feature} & \textbf{Importance} \\
		\hline
		CustomerType\_Occasionnel & 18.7\% \\
		Frequency & 15.6\% \\
		FirstPurchaseDaysAgo & 14.5\% \\
		UniqueInvoices & 11.5\% \\
		CustomerType\_Régulier & 6.5\% \\
		\hline
	\end{tabular}
\end{table}
\subsection{Interprétation}
Les clients "Occasionnels" et la fréquence d'achat sont les meilleurs prédicteurs du churn. Un client qui achète rarement et depuis longtemps a un risque de départ significativement plus élevé.
\section{Régression — Prédiction de MonetaryTotal}

\subsection{Objectif}
Prédire le montant total dépensé par un client (valeur continue). Features : 42 colonnes numériques.

\subsection{Colonnes Exclues}
\texttt{MonetaryAvg}, \texttt{MonetaryStd}, \texttt{MonetaryMin}, \texttt{MonetaryMax} (toutes dérivées de \texttt{MonetaryTotal} → leakage de régression).

\subsection{Résultats}

\begin{table}[H]
	\centering
	\caption{Performances des modèles de régression}
	\label{tab:regression}
	\begin{tabular}{|l|c|c|c|}
		\hline
		\textbf{Modèle} & \textbf{MAE (£)} & \textbf{RMSE (£)} & \textbf{R²} \\
		\hline
		Linear Regression & 819.82 & 6584.67 & 0.6630 \\
		\textbf{Random Forest} & \textbf{688.07} & \textbf{6361.77} & \textbf{0.6855} \\
		XGBoost & 825.50 & 7996.71 & 0.5030 \\
		\hline
	\end{tabular}
\end{table}
\begin{figure}[H]
	\centering
	\includegraphics[width=0.8\textwidth]{images/regression_pred_vs_real.png}
	\caption{Régression Random Forest — Valeurs réelles vs prédites. La droite rouge représente la prédiction parfaite. Les points s'éloignant de la diagonale correspondent aux clients VIP.}
	\label{fig:regression}
\end{figure}
\subsection{Interprétation des Résultats}
\begin{itemize}
	\item \textbf{R² = 0.6855} : le modèle explique 68.55\% de la variance du montant dépensé — performance correcte pour ce type de données
	\item \textbf{RMSE élevé (6361£)} : s'explique par la présence de clients VIP (Cluster 3) qui dépensent des centaines de milliers de livres, créant des valeurs extrêmes difficiles à prédire précisément
	\item \textbf{MAE = 688£} : l'erreur médiane est plus raisonnable, confirmant que le RMSE est tiré vers le haut par quelques outliers
\end{itemize}
\section{Prédiction et Déploiement des Modèles}

\subsection{Architecture du Script \texttt{predict.py}}

Le script \texttt{predict.py} constitue l'interface de production des modèles entraînés. Il charge les trois artefacts sauvegardés et expose des fonctions de prédiction réutilisables pour une intégration dans une application web ou un pipeline décisionnel.

\begin{table}[H]
	\centering
	\caption{Modèles chargés par \texttt{predict.py}}
	\label{tab:models_loaded}
	\begin{tabular}{|l|l|l|}
		\hline
		\textbf{Fichier} & \textbf{Modèle} & \textbf{Type} \\
		\hline
		\texttt{models/best\_classifier.pkl} & XGBClassifier & Classification (Churn) \\
		\texttt{models/best\_regressor.pkl} & RandomForestRegressor & Régression (MonetaryTotal) \\
		\texttt{models/kmeans\_model.pkl} & KMeans (k=4) & Clustering (Segmentation) \\
		\hline
	\end{tabular}
\end{table}
\begin{table}[H]
	\centering
	\caption{Pipeline d'exécution de \texttt{predict.py}}
	\label{tab:pipeline_predict}
	\begin{tabular}{|l|l|l|l|}
		\hline
		\textbf{Étape} & \textbf{Description} & \textbf{Entrée} & \textbf{Sortie} \\
		\hline
		1 & Chargement des modèles & \texttt{models/*.pkl} & 3 objets Python \\
		2 & Récupération colonnes référence & \texttt{X\_train.csv} & Listes de colonnes \\
		3 & Fonctions de prédiction & DataFrame client & Prédictions \\
		4 & Démonstration sur test & \texttt{X\_test.csv} (875) & Rapport complet \\
		5 & Génération rapport combiné & Toutes prédictions & \texttt{predictions.csv} \\
		6 & Prédiction unitaire & Valeurs normalisées & Résultat individuel \\
		\hline
	\end{tabular}
\end{table}
\begin{table}[H]
	\centering
	\caption{Modèles chargés par \texttt{predict.py}}
	\label{tab:models_loaded}
	\begin{tabular}{|l|l|l|}
		\hline
		\textbf{Fichier} & \textbf{Modèle} & \textbf{Type} \\
		\hline
		\texttt{models/best\_classifier.pkl} & XGBClassifier & Classification (Churn) \\
		\texttt{models/best\_regressor.pkl} & RandomForestRegressor & Régression (MonetaryTotal) \\
		\texttt{models/kmeans\_model.pkl} & KMeans (k=4) & Clustering (Segmentation) \\
		\hline
	\end{tabular}
\end{table}
\begin{table}[H]
	\centering
	\caption{Pipeline d'exécution de \texttt{predict.py}}
	\label{tab:pipeline_predict}
	\begin{tabular}{|l|l|l|l|}
		\hline
		\textbf{Étape} & \textbf{Description} & \textbf{Entrée} & \textbf{Sortie} \\
		\hline
		1 & Chargement des modèles & \texttt{models/*.pkl} & 3 objets Python \\
		2 & Récupération colonnes référence & \texttt{X\_train.csv} & Listes de colonnes \\
		3 & Fonctions de prédiction & DataFrame client & Prédictions \\
		4 & Démonstration sur test & \texttt{X\_test.csv} (875) & Rapport complet \\
		5 & Génération rapport combiné & Toutes prédictions & \texttt{predictions.csv} \\
		6 & Prédiction unitaire & Valeurs normalisées & Résultat individuel \\
		\hline
	\end{tabular}
\end{table}
\subsubsection{Seuils de Risque Churn}

La fonction \texttt{predict\_churn} catégorise la probabilité de départ selon quatre niveaux de risque opérationnels :

\begin{table}[H]
	\centering
	\caption{Seuils de risque Churn}
	\label{tab:risk_thresholds}
	\begin{tabular}{|l|c|l|}
		\hline
		\textbf{Niveau} & \textbf{Seuil probabilité} & \textbf{Signification} \\
		\hline
		Faible & $< 0.25$ & Client stable, peu de risque \\
		Moyen & $0.25 - 0.50$ & Surveillance recommandée \\
		Élevé & $0.50 - 0.75$ & Intervention conseillée \\
		Critique & $> 0.75$ & Action immédiate requise \\
		\hline
	\end{tabular}
\end{table}
\begin{tcolorbox}[colback=codegreen!10, colframe=codegreen, title=\textbf{Note sur \texttt{predict\_monetary}}]
	Les valeurs négatives prédites sont clippées à 0 via \texttt{np.maximum(predictions, 0)}. Dans le jeu de données d'entraînement, les montants négatifs correspondent à des remboursements historiques, phénomène non prédictible pour un futur client sans historique d'achat.
\end{tcolorbox}
\subsection{Résultats sur le Jeu de Test (875 Clients)}

\subsubsection{Résultats de la Classification}

Le modèle XGBoost atteint une accuracy de 99.77\% sur l'ensemble de test.

\begin{table}[H]
	\centering
	\caption{Distribution des prédictions de churn}
	\label{tab:pred_distribution}
	\begin{tabular}{|l|c|c|c|}
		\hline
		\textbf{Classe} & \textbf{Libellé} & \textbf{Nombre} & \textbf{Pourcentage} \\
		\hline
		0 & Fidèle prédit & 586 & 67.0\% \\
		1 & Churné prédit & 289 & 33.0\% \\
		\hline
		\textbf{Total} & & \textbf{875} & \textbf{100\%} \\
		\hline
	\end{tabular}
\end{table}
\begin{table}[H]
	\centering
	\caption{Distribution des niveaux de risque}
	\label{tab:risk_distribution}
	\begin{tabular}{|l|c|c|}
		\hline
		\textbf{Niveau} & \textbf{Nombre} & \textbf{Pourcentage} \\
		\hline
		Faible & 585 & 66.9\% \\
		Moyen & 1 & 0.1\% \\
		Critique & 289 & 33.0\% \\
		\hline
	\end{tabular}
\end{table}
\begin{figure}[H]
	\centering
	\includegraphics[width=0.7\textwidth]{reports/predict_risk_distribution.png}
	\caption{Distribution des niveaux de risque Churn sur les 875 clients de test}
	\label{fig:risk_distribution}
\end{figure}
\begin{figure}[H]
	\centering
	\includegraphics[width=0.7\textwidth]{reports/predict_proba_distribution.png}
	\caption{Distribution des probabilités de churn — concentration aux extrêmes (nature synthétique du dataset)}
	\label{fig:proba_distribution}
\end{figure}
\begin{tcolorbox}[colback=warningred!10, colframe=warningred, title=\textbf{Observation — Distribution Binaire des Risques}]
	Le modèle est très binaire : 99.9\% des clients sont classés en niveau \textit{Faible} ou \textit{Critique}, avec un seul client en zone \textit{Moyen}. Cette concentration aux extrêmes confirme la nature synthétique du dataset, où le Churn a été généré de façon déterministe par une formule mathématique. Dans un contexte réel de e-commerce, la distribution serait plus étalée sur les quatre niveaux de risque, avec davantage de cas ambigus.
\end{tcolorbox}
\subsubsection{Résultats de la Segmentation}

Le clustering K-Means identifie quatre segments clients cohérents avec les profils RFM.

\begin{table}[H]
	\centering
	\caption{Distribution des segments clients}
	\label{tab:segment_distribution}
	\begin{tabular}{|l|c|c|l|}
		\hline
		\textbf{Segment} & \textbf{Nombre} & \textbf{Pourcentage} & \textbf{Profil} \\
		\hline
		Réguliers & 642 & 73.4\% & Comportement moyen équilibré \\
		Dormants & 221 & 25.3\% & Inactifs, risque churn élevé \\
		Champions & 11 & 1.3\% & Très actifs, dépenses élevées \\
		VIP & 1 & 0.1\% & Ultra-premium \\
		\hline
	\end{tabular}
\end{table}
\begin{figure}[H]
	\centering
	\includegraphics[width=0.7\textwidth]{reports/predict_segments.png}
	\caption{Distribution des segments clients K-Means sur les 875 clients de test}
	\label{fig:segments}
\end{figure}
\subsubsection{Résultats de la Régression}

Le modèle Random Forest Regressor prédit le montant total dépensé avec un R² de 0.6855 sur l'ensemble d'entraînement.

\begin{table}[H]
	\centering
	\caption{Statistiques des montants prédits}
	\label{tab:monetary_stats}
	\begin{tabular}{|l|c|}
		\hline
		\textbf{Statistique} & \textbf{Valeur} \\
		\hline
		Minimum & 0.00 £ \\
		Médiane & 0.00 £ \\
		Moyenne & 153.60 £ \\
		Maximum & 1645.93 £ \\
		\hline
	\end{tabular}
\end{table}
\begin{figure}[H]
	\centering
	\includegraphics[width=0.8\textwidth]{reports/predict_monetary_by_segment.png}
	\caption{Distribution des montants prédits par segment (clients avec montant > 0£)}
	\label{fig:monetary_by_segment}
\end{figure}
\begin{tcolorbox}[colback=lightgray, colframe=accentorange, title=\textbf{Observation — Médiane à 0£}]
	La médiane à 0.00£ indique que plus de 50\% des clients ont un montant prédit très faible (clippé à zéro). Le modèle Random Forest (R²=0.6855) est plus précis sur les gros montants (segments Champions et VIP) que sur les petits clients. Cette caractéristique est cohérente avec la nature des données : les clients à faible activité sont plus difficiles à modéliser.
\end{tcolorbox}
\subsection{Comparaison Churn Réel vs Prédit}

\begin{figure}[H]
	\centering
	\includegraphics[width=0.8\textwidth]{reports/predict_real_vs_pred.png}
	\caption{Comparaison churn réel vs prédit — distributions quasi-identiques confirmant l'accuracy de 99.77\%}
	\label{fig:real_vs_pred}
\end{figure}

Les diagrammes circulaires montrent une quasi-identité entre la distribution réelle du churn et celle prédite par le modèle, avec respectivement 33.3\% de churnés dans les données réelles et 33.0\% dans les prédictions.
\subsection{Rapport Combiné — Extrait \texttt{predictions.csv}}

Le script génère un fichier \texttt{predictions.csv} contenant l'ensemble des prédictions pour les 875 clients de test.

\begin{table}[H]
	\centering
	\small
	\caption{Extrait des prédictions (10 premiers clients)}
	\label{tab:predictions_extract}
	\begin{tabular}{|c|c|c|c|c|c|c|}
		\hline
		\textbf{Client} & \textbf{Churn Réel} & \textbf{Churn Prédit} & \textbf{Probabilité} & \textbf{Risque} & \textbf{Montant £} & \textbf{Segment} \\
		\hline
		0 & 0 & 0 & 0.22\% & Faible & 0.00 & Réguliers \\
		1 & 0 & 0 & 0.14\% & Faible & 0.00 & Réguliers \\
		2 & 1 & 1 & 99.94\% & Critique & 0.00 & Dormants \\
		3 & 0 & 0 & 0.21\% & Faible & 0.00 & Réguliers \\
		4 & 1 & 1 & 96.83\% & Critique & 853.02 & Réguliers \\
		5 & 0 & 0 & 0.04\% & Faible & 0.00 & Réguliers \\
		6 & 1 & 1 & 99.97\% & Critique & 0.00 & Dormants \\
		7 & 1 & 1 & 99.95\% & Critique & 0.00 & Dormants \\
		8 & 0 & 0 & 1.77\% & Faible & 853.02 & Réguliers \\
		9 & 0 & 0 & 0.21\% & Faible & 851.55 & Réguliers \\
		\hline
		\end{tabular}
	\end{table}
\begin{tcolorbox}[colback=backcolour, colframe=headerblue, title=\textbf{Observations sur l'extrait}]
	\begin{itemize}
		\item \textbf{Client 2} : Dormant avec 99.94\% de probabilité de churn — cohérence parfaite entre segmentation et classification.
		\item \textbf{Client 4} : Churné mais classé Régulier — un client au comportement régulier peut néanmoins partir, illustrant que la fidélité passée ne garantit pas la fidélité future.
		\item \textbf{Clients 8-9} : Fidèles avec 853£ de montant prédit — clients rentables à conserver absolument, cibles prioritaires pour des actions de fidélisation.
	\end{itemize}
\end{tcolorbox}
\subsection{Exemple de Prédiction pour un Nouveau Client}

Le script permet également de prédire pour un client fictif dont les caractéristiques ont été normalisées avec le \texttt{StandardScaler} sauvegardé.

\begin{table}[H]
	\centering
	\caption{Caractéristiques du nouveau client (valeurs normalisées)}
	\label{tab:new_customer_features}
	\begin{tabular}{|l|c|}
		\hline
		\textbf{Feature} & \textbf{Valeur normalisée} \\
		\hline
		Recency & -0.4 (actif récemment) \\
		Frequency & 0.2 \\
		FirstPurchaseDaysAgo & -0.3 (client récent) \\
		\hline
	\end{tabular}
\end{table}
\begin{table}[H]
	\centering
	\caption{Résultats de prédiction pour le nouveau client}
	\label{tab:new_customer_results}
	\begin{tabular}{|l|c|}
		\hline
		\textbf{Indicateur} & \textbf{Valeur} \\
		\hline
		Churn prédit & Churné (1) \\
		Probabilité & 90.6\% \\
		Niveau de risque & Critique \\
		Montant prédit & 3.92 £ \\
		Segment & Réguliers \\
		\hline
	\end{tabular}
\end{table}
\begin{tcolorbox}[colback=warningred!10, colframe=warningred, title=\textbf{Analyse — Nouveau Client à Risque Critique}]
	Malgré une activité récente (\texttt{Recency} = -0.4, soit inférieur à la moyenne) et un premier achat récent (\texttt{FirstPurchaseDaysAgo} = -0.3), le modèle prédit un churn critique à 90.6\%. Cette prédiction s'explique par l'importance des features \texttt{CustomerType\_Occasionnel} (18.7\%) et \texttt{Frequency} (15.6\%) dans le modèle XGBoost. Un client avec peu de factures (\texttt{Frequency} légèrement positive) et sans historique de fidélisation est systématiquement perçu comme à risque de départ, même s'il vient d'effectuer un achat.
\end{tcolorbox}
\subsection{Limites et Recommandations pour le Déploiement}

\subsubsection{Limites Identifiées}

\begin{table}[H]
	\centering
	\caption{Limites du système de prédiction}
	\label{tab:limitations}
	\begin{tabular}{|l|l|l|}
		\hline
		\textbf{Limite} & \textbf{Cause} & \textbf{Impact} \\
		\hline
		Prédictions très binaires & Dataset synthétique & Peu de zones grises, décisions binaires \\
		Médiane régression = 0£ & Sous-estimation petits clients & Imprécision pour clients modestes \\
		Valeurs normalisées requises & StandardScaler appliqué & Preprocessing obligatoire en production \\
		\hline
	\end{tabular}
\end{table}
\subsubsection{Recommandations pour le Déploiement en Production}


\end{document}